\documentclass[12pt, titlepage]{article}
\pdfinfoomitdate=1
\pdftrailerid{}

\usepackage{booktabs}
\usepackage{tabularx}
\usepackage{hyperref}
\hypersetup{
    colorlinks,
    citecolor=black,
    filecolor=black,
    linkcolor=red,
    urlcolor=blue
}
\usepackage[round]{natbib}

\input{../Comments}
\input{../Common}

% Links to project documents (relative to docs/VnVReport/)
\newcommand{\SRS}{\href{../SRS/SRS.pdf}{SRS}}
\newcommand{\VnVPlan}{\href{../VnVPlan/VnVPlan.pdf}{VnV Plan}}


\begin{document}

\title{Verification and Validation Report: \progname}
\author{\authname}
\date{}

\maketitle

\pagenumbering{roman}

\section{Revision History}

\begin{tabularx}{\textwidth}{p{3cm}p{2cm}X}
  \toprule {\bf Date} & {\bf Version} & {\bf Notes} \\
  \midrule
  2026-03-04          & 1.0           & Add         \\
  Date 2              & 1.1           & Notes       \\
  \bottomrule
\end{tabularx}

~\newpage

\section{Symbols, Abbreviations and Acronyms}

\renewcommand{\arraystretch}{1.2}
\begin{tabular}{l l}
  \toprule
  \textbf{symbol} & \textbf{description} \\
  \midrule
  T               & Test                 \\
  \bottomrule
\end{tabular}\\

\wss{symbols, abbreviations or acronyms -- you can reference the SRS tables if needed}

\newpage

\tableofcontents

\listoftables %if appropriate

\listoffigures %if appropriate

\newpage

\pagenumbering{arabic}

This document ...

\section{Functional Requirements Evaluation}
\label{sec:fr_eval}

This section reports the verification evidence for \textbf{functional
  requirements} defined in the \SRS{}. Test procedures and test case definitions
are specified in the \VnVPlan{}; to avoid unnecessary repetition, this report
focuses on \textbf{execution status}, \textbf{results}, and \textbf{evidence}.

\subsection{Scope and Execution Policy}
\label{sec:fr_scope}

As stated in the latest \VnVPlan{} (Revision 5.0), the only testing executed
for this project is \textbf{unit testing} and a \textbf{stakeholder review}
(McMaster Rocketry Team). Integration/system/field tests are documented in the
plan for traceability but are \textbf{not executed} due to resource and contact
constraints.

\subsection{Summary of Planned Functional Tests vs Execution}
\label{sec:fr_summary_table}

Table~\ref{tab:fr_exec_summary} maps the planned functional test IDs (from the
\VnVPlan{}) to (1) their related SRS functional requirements and (2) the actual
execution status in this project. Where tests are executed, evidence is linked
in the final column.

\begin{table}[H]
  \centering
  \caption{Functional Test Execution Summary (traceable to \VnVPlan{})}
  \label{tab:fr_exec_summary}
  \renewcommand{\arraystretch}{1.2}
  \begin{tabularx}{\textwidth}{p{3.4cm} p{3.2cm} p{3.0cm} X}
    \toprule
    \textbf{Test ID (VnVPlan)} & \textbf{FR(s) (SRS)} & \textbf{Execution} & \textbf{Evidence / Notes}   \\
    \midrule

    test-aquire-frame          & FR-1                 & Not executed       & Out of scope (system test). \\
    test-disconnect-camera     & FR-1                 & Not executed       & Out of scope (system test). \\
    test-hdmi-output           & FR-8, FR-9           & Not executed       & Out of scope (system test). \\
    test-preview-real-time     & FR-10                & Not executed       & Out of scope (system test). \\
    test-network-disconnect    & FR-10                & Not executed       & Out of scope (system test). \\

    test-recording             & FR-11                & Not executed       & Out of scope (system test). \\
    test-disk-full             & FR-11                & Not executed       & Out of scope (system test). \\
    test-manage-recordings     & FR-12                & Not executed       & Out of scope (system test). \\

    test-manual-control        & FR-2, FR-6           & Not executed       & Out of scope (system test).
    \\ test-transition-to-tracking & FR-3, FR-4 & Not executed & Out of scope
    (system test).                                                                                       \\ test-tracking & FR-2, FR-3, FR-7 & Not executed & Out of
    scope (system test).                                                                                 \\ test-return-to-idle & FR-5 & Not executed & Out of
    scope (system test).                                                                                 \\ test-exit-gimbal-range & FR-2, FR-7 & Not executed &
    Out of scope (system test).                                                                          \\

    \bottomrule
  \end{tabularx}
\end{table}

\noindent
\textbf{Note:} Functional requirements are primarily verified through unit tests
in this project. The unit-level test suite (names, coverage, and results) is
reported in Section~\ref{sec:unit_testing} and traced back to requirements in
Section~\ref{sec:trace_reqs}. The above table is retained to preserve
traceability to the \VnVPlan{} test IDs even when those tests are not executed.

\subsection{Unit-Test Evidence for Functional Requirements}
\label{sec:fr_unit_evidence}

Because system-level FR tests were not executed, FR verification relies on unit
tests that validate the underlying behavior (e.g., state transitions, command
validation, recording pipeline logic, UI action handlers). Evidence is
presented in two parts:

\begin{enumerate}
  \item \textbf{Requirement-to-test mapping:} Section~\ref{sec:trace_reqs} lists
        which unit tests support each FR.
  \item \textbf{Execution results:} Section~\ref{sec:unit_testing} reports
        pass/fail status and, where applicable, links to CI logs or local test output.
\end{enumerate}

\subsubsection{Evidence Format and Reproducibility}
For each functional requirement verified by unit testing, the report records:
\begin{itemize}
  \item The FR identifier(s) from the \SRS{}.
  \item The unit test file(s) / test case names (self-descriptive naming).
  \item The execution outcome (pass/fail) for the final submission revision.
  \item A pointer to evidence (CI job log, screenshot, or attached test report).
\end{itemize}

\subsection{Stakeholder Validation Touchpoint}
\label{sec:fr_stakeholder_touchpoint}

Stakeholder validation is performed via a review session with the McMaster
Rocketry Team as defined in the \VnVPlan{}. The outcomes of this review (action
items, approvals, and required changes) are summarized in
Section~\ref{sec:changes_due_testing}. Any requested changes are tracked as
GitHub Issues and linked there to maintain traceability.

\section{Nonfunctional Requirements Evaluation}

\subsection{Usability}

\subsection{Performance}

\subsection{etc.}

\section{Comparison to Existing Implementation}

This section will not be appropriate for every project.

\section{Unit Testing}

\section{Changes Due to Testing}

\wss{This section should highlight how feedback from the users and from
  the supervisor (when one exists) shaped the final product.  In particular
  the feedback from the Rev 0 demo to the supervisor (or to potential users)
  should be highlighted.}

\section{Automated Testing}

\section{Trace to Requirements}

\section{Trace to Modules}

\section{Code Coverage Metrics}

\bibliographystyle{plainnat}
\bibliography{../../refs/References}

\newpage{}
\section*{Appendix --- Reflection}

The information in this section will be used to evaluate the team members on
the graduate attribute of Reflection.

\input{../Reflection.tex}

\begin{enumerate}
  \item What went well while writing this deliverable?
  \item What pain points did you experience during this deliverable, and how did you
        resolve them?
  \item Which parts of this document stemmed from speaking to your client(s) or a proxy
        (e.g. your peers)? Which ones were not, and why?
  \item In what ways was the Verification and Validation (VnV) Plan different from the
        activities that were actually conducted for VnV? If there were differences,
        what changes required the modification in the plan? Why did these changes
        occur? Would you be able to anticipate these changes in future projects? If
        there weren't any differences, how was your team able to clearly predict a
        feasible amount of effort and the right tasks needed to build the evidence that
        demonstrates the required quality? (It is expected that most teams will have
        had to deviate from their original VnV Plan.)
\end{enumerate}

\end{document}